% !PRINTED 612
\weberpage{612}

%<<CONT
then

\[ \chi = \frac{Q\varphi}{\xi} = \beta_1t_1+\beta_2t_2, \]

and consequently $\beta_1,\beta_2$ are integers divisible by $\psi'$.

If we form from (9) the discriminant $\varDelta(\beta_1,\beta_2)$, we obtain

\weq{10}{\varDelta(\beta_1,\beta_2) = \begin{vmatrix}\beta_1, & \beta_2\\ \beta_1', & \beta_2'\end{vmatrix}^2 = \frac{e^4c^4Q^4}{(\xi\xi')^2}(\omega-\omega')^2.}

But from (5) and (7) it follows that

\[ \xi\xi' = N(\xi) = PQ = e^2cQ, \]

and if we further set $\omega-\omega' = \dfrac{\sqrt{\varDelta}}{c}$, we find

\weq{11}{\varDelta(\beta_1,\beta_2) = Q^2\varDelta.}

But from this it follows, by the theorem of \S.\,147, 2., that $\beta_1,\beta_2$ is a basis of $\psi'$.

On the other hand, by (6), $e'c', e'c'\eta'$ is also a basis of $\psi'$, and consequently the integral rational numbers $p,q,r,s$ can be so determined that

\weq{12}{\begin{aligned}
\beta_1 &= \frac{Qec}{\xi} = e'c'(r\eta'+s)\\
\beta_2 &= \frac{Qec}{\xi}\,\omega = e'c'(p\eta'+q),
\end{aligned}}

and at the same time

\weq{13}{ps-qr = \pm 1.}

But from (12) it follows by division:

\weq{14}{\omega = \frac{p\eta'+q}{r\eta'+s},}

that is, $\omega$ is equivalent to $\eta'$.

If we now replace $\varphi$ by the equivalent form $\varphi_1$, there subsists an equation like (5):

\[ \varphi_1\psi = \varepsilon_1\xi_1, \]

in which $\psi$ has remained unchanged. Consequently $\omega_1$ is also equivalent to $\eta$, and hence $\omega$ and $\omega_1$ are equivalent to one another. % [sic] source prints $\eta$; repeating the derivation with unchanged $\psi$ requires $\eta'$
But this is the assertion expressed in our Theorem 3., which is therewith proved.

According to this, then, to every ideal class of the field $\Omega$ there corresponds a class of equivalent numbers $\omega$, and conversely; and we may therefore add the further theorem:
%CONT>>
